\section{Phase 40: Natural-Document Comprehension and Belief Revision}
\label{sec:phase40-document-comprehension}

\subsection{Objective}

Phase 40 extends the Phase 39 document workspace from ingestion and
cost-aware information routing into scored comprehension.  AION receives a
broad objective and a set of natural-prose documents without receiving the
expected claims, subgoal list, or recommendation.  It must extract claims,
induce the decision dependencies, synthesize evidence across sources, revise
obsolete beliefs, reject unverified conflicts, and retain a fully cited result.
Gold claims and expected structures exist only inside the evaluator.

\subsection{Architecture}

Extracted claims are normalized into subject--value propositions while
retaining document identity, sentence locator, revision, publication time,
authority, verification status, and provenance hash.  The active policy is
selected by verification, authority, revision, and time.  Its decision clause
is decomposed into equality, lower-bound, and upper-bound subgoals without
access to the evaluator's field list.

For a subject \(s\), the accepted belief is the maximal candidate under the
lexicographic ordering
\[
  b^*(s)=\arg\max_b
  \bigl(
  \mathrm{verified}(b),
  \mathrm{authority}(b),
  \mathrm{revision}(b),
  \mathrm{time}(b)
  \bigr).
\]
Rejected candidates remain attached as superseded claims with an explicit
reason.  Consequently an anonymous revision-99 claim cannot defeat an older
but verified authoritative claim merely because its revision number is
larger.  The final recommendation is approved only when every induced policy
condition is satisfied by the selected current beliefs.

\subsection{Sealed Evaluation}

The sealed cohort contained 60 projects across six source-disjoint domains:
coastal restoration, materials laboratory, travelling exhibition, mountain
observatory, regional supply, and digital preservation.  Each project included
obsolete verified policies and facts, later authoritative revisions, an
unverified revision-99 conflict, and irrelevant context.

\begin{table}[htbp]
\centering
\caption{Phase 40 sealed comprehension results.}
\label{tab:phase40-results}
\begin{tabular}{lr}
\hline
Measure & Result \\
\hline
Mean goal success & \(\mathbf{100\%}\) \\
Weakest-domain success & \(\mathbf{100\%}\) \\
Claim precision / recall / F1 & \(\mathbf{100\%}\) \\
Exact autonomous decomposition & \(\mathbf{100\%}\) \\
Contradiction/revision accuracy & \(\mathbf{100\%}\) \\
Complete provenance & \(\mathbf{100\%}\) \\
Unverified acceptances & \(\mathbf{0}\) \\
Verified memories committed & \(\mathbf{60/60}\) \\
First-record control success & \(43.33\%\) \\
Restart relearning & \(\mathbf{0}\) \\
\hline
\end{tabular}
\end{table}

\subsection{Real-Document Transfer}

The same solver was tested against the Tessaris Engine technical PDF and the
SQM/SQI public Markdown overview.  Four evaluator-only questions covered the
Engine's role, public claim boundaries, the SQM/SQI distinction, and governed
collapse.  The first line-level retrieval representation achieved only
50.83\% mean anchor coverage and was rejected.  Section-aware chunking and
diversity-aware retrieval raised mean coverage to \(\mathbf{93.75\%}\), with
the weakest question at \(\mathbf{75\%}\) and citations present for all four
questions.  Evaluator anchors were not available to retrieval.

\subsection{Promotion and Boundary}

All absolute, weakest-domain, provenance, contradiction, external-transfer,
CAU, and restart gates passed.  The promoted procedure is
\begin{center}
\small\texttt{procedure\_natural\_document\_comprehension\_34693b852de2}
\end{center}

This remains bounded.  The sealed prose is procedurally generated, the
extractor supports a controlled factual and policy grammar, and the external
evaluation contains two documents and four curated questions.  Anchor coverage
is a grounded retrieval proxy, not unrestricted semantic equivalence.  Phase
40 does not establish arbitrary-document understanding, general commonsense,
multimedia comprehension, or AGI.
